% ---------------------------------------------------------------------
%  PHỤ LỤC
% ---------------------------------------------------------------------
% Trong \backmatter thì \chapter đã tự thêm vào mục lục, không cần
% \addcontentsline nữa, thêm vào là trùng neo dấu trang.
\chapter{Phụ lục}
\markboth{Phụ lục}{}

\begin{chuy}
\goiy{phụ lục là nơi đặt những nội dung dài không tiện chèn vào thân báo
cáo: mã nguồn đầy đủ, bảng dữ liệu mở rộng, kết quả thử nghiệm bổ sung, biên
bản hướng dẫn. Riêng mục A và mục B nên giữ lại.}
\end{chuy}

\section*{A. Hướng dẫn cài đặt và chạy chương trình}
\phantomsection
\addcontentsline{toc}{section}{A. Hướng dẫn cài đặt và chạy chương trình}

Viết các bước để một người chưa từng thấy đồ án vẫn chạy được sản phẩm. Hội
đồng sẽ chạy thử theo đúng các bước này.

\begin{lstlisting}[language=bash, caption={Các bước cài đặt và khởi chạy ứng dụng.}]
# 1. Clone the repository
git clone https://github.com/tencuaem/ten-do-an.git
cd ten-do-an

# 2. Create a virtual environment and install dependencies
python -m venv .venv
.venv\Scripts\activate
pip install -r requirements.txt

# 3. Train the model (skip if a trained checkpoint is provided)
python train.py --epochs 30 --data data/

# 4. Launch the web application, then open http://localhost:5000
python app.py
\end{lstlisting}

\section*{B. Bảng tự kiểm trước khi nộp}
\phantomsection
\addcontentsline{toc}{section}{B. Bảng tự kiểm trước khi nộp}

Trước khi in quyển, em tự đánh dấu vào từng ô dưới đây.

\begin{tukiem}
  \item Sản phẩm chạy được trên máy khác, không chỉ trên máy của em.
  \item Đã chạy thử demo trên đúng môi trường sẽ dùng khi bảo vệ.
  \item Mã nguồn nằm trên kho \code{git}, lịch sử commit cho thấy tiến độ
        trải đều chứ không dồn vào tuần cuối.
  \item Có tệp \code{README} và \code{requirements.txt} để người khác chạy lại.
  \item Mọi hình, bảng và giải thuật đều có chú thích, có nhãn, và được nhắc
        tới trong thân bài.
  \item Mọi mục tiêu ở Chương~\ref{ch:modau} đều được đối chiếu ở
        Chương~\ref{ch:ketluan}.
  \item Mọi con số trong báo cáo trùng khớp với kết quả chạy thật, không có
        số nào là ước lượng hay bịa.
  \item Mọi tài liệu trong \code{refs.bib} đều được trích dẫn ít nhất một
        lần, và mọi nguồn dữ liệu, mã tham khảo đều đã ghi rõ.
  \item Đã khai báo phần nào có dùng công cụ trí tuệ nhân tạo hỗ trợ, và em
        giải thích được toàn bộ hệ thống khi bị hỏi.
  \item Đã xoá hết hộp gợi ý, mọi khối nội dung mẫu, và cả
        Chương~\ref{ch:soanthao} vốn chỉ là sổ tra cứu.
\end{tukiem}

\begin{luuy}
Hai mục cuối là lỗi hay gặp nhất khi dùng mẫu có sẵn. Trước khi nộp hãy tìm
trong bản PDF xem còn dòng chữ xám \emph{NỘI DUNG MẪU} nào sót lại không.
\end{luuy}

\section*{C. Nhận xét của giảng viên hướng dẫn}
\phantomsection
\addcontentsline{toc}{section}{C. Nhận xét của giảng viên hướng dẫn}

\khungnhanxet[7.5cm]

\vspace{0.8em}
\noindent\begin{minipage}[t]{0.5\textwidth}
\color{daunavy}
{\color{daumaroon}\bfseries Điểm số}

\vspace{0.7em}
\dongchamnhan{Bằng số:}

\vspace{0.9em}
\dongchamnhan{Bằng chữ:}
\end{minipage}\hfill
\okyten{\dauGVHD}{Giảng viên hướng dẫn}
